(1) Itô公式 \(d(x^2)=2x\,dx+(dx)^2\) 给出
\[
\dot{\langle x\rangle}=-(k/\zeta)\langle x\rangle,\quad
\dot{\langle x^2\rangle}=-(2k/\zeta)\langle x^2\rangle+2k_BT_f/\zeta.
\]
初态均值0、方差 \(\sigma_i^2=k_BT_i/k\)。过程线性且保持高斯性，故
\[
p(x,t)=\frac{e^{-x^2/(2\sigma_t^2)}}{\sqrt{2\pi\sigma_t^2}},\quad
\boxed{\sigma_t^2=\sigma_f^2+(\sigma_i^2-\sigma_f^2)e^{-2kt/\zeta}},
\quad\sigma_f^2=k_BT_f/k.
\]
方差弛豫时间 \(\zeta/(2k)\)，不是均值弛豫时间 \(\zeta/k\)。

(2) 平均能量 \(U=k\langle x^2\rangle/2\)，故
\(U(t)=k_B[T_f+(T_i-T_f)e^{-2kt/\zeta}]/2\)。
谐阱参数不随时间改变，外功 \(W=0\)；第一定律给进入系统的热量
\[
\boxed{Q(t)=\Delta U=\tfrac12k_B(T_f-T_i)(1-e^{-2kt/\zeta})}.
\]
当冷却 \(T_f<T_i\) 时 \(Q<0\) 表示系统放热，最终 \(U=k_BT_f/2\)，满足单个位置自由度的能量均分。

(3) 系统熵 \(S_{\rm sys}(t)=\tfrac{k_B}{2}\ln(2\pi e\sigma_t^2)\)，热浴熵 \(\Delta S_{\rm bath}=-Q/T_f\)。令 \(r_i=T_i/T_f\)，\(r_t=\sigma_t^2/\sigma_f^2=1+(r_i-1)e^{-2kt/\zeta}\)，则
\[
\boxed{\Delta S_{\rm tot}(t)=\frac{k_B}{2}[\ln(r_t/r_i)+r_i-r_t]}.
\]
求导得 \(\dot S_{\rm tot}=(k_Bk/\zeta)(r_t-1)^2/r_t\ge0\)。由于初始为0，总熵产生任意时刻非负；终态为 \(\tfrac{k_B}{2}(r_i-1-\ln r_i)\ge0\)，后式不可误当任意时刻公式。

(4) 相对最终平衡高斯态的KL散度
\(D(p_t\Vert p_f)=\tfrac12(r_t-1-\ln r_t)\)，其导数
\[
\dot D=-\frac{k}{\zeta}\frac{(r_t-1)^2}{r_t}\le0.
\]
非平衡自由能差 \(F[p_t]-F[p_f]=k_BT_fD\) 单调下降，且 \(-\dot F=T_f\dot S_{\rm tot}\)，独立核对了熵产生率。若 \(T_i=T_f\) 所有变化均为零；严格 \(T_f=0\) 时浴熵公式含 \(1/T_f\) 且高斯平衡分布退化，必须先在正温范围内推导。
